\section{Theory-to-Lean Correspondence}

Traceability is reported at theorem granularity, mapping named theory claims to
Lean theorem identifiers.

\subsection{Status Semantics}
\label{sec:status-semantics}

We use the following labels in correspondence tables:
\begin{itemize}
  \item \textbf{Proved}: theorem stated in theory text and proved in Lean.
  \item \textbf{Demoted}: originally an axiom, now proved in Lean.
  \item \textbf{Novel}: theorem emerged during formalization and was not explicitly stated in the original theory text.
  \item \textbf{Lean-only}: proved in Lean but not stated in the Theory Paper (auxiliary lemmas, witnesses, or obligations surfaced during proof construction).
\end{itemize}

\subsection{Coverage and Mapping Strategy}

Mappings prioritize exact theorem names and principal Lean counterparts.
Multi-lemma realizations are represented by a principal theorem name plus
supporting theorem names where required for faithful traceability.

\subsection{Chapter-Wise Mapping Policy}

\begin{enumerate}
  \item Claims presented as theorems in theory chapters are labeled \textbf{Proved}.
  \item Former theory axioms shown derivable in Lean are labeled \textbf{Demoted}.
  \item New theorems surfaced from Lean proof obligations are labeled \textbf{Novel}.
\end{enumerate}

This policy avoids conflating routine theorem verification with axiom demotion,
and preserves a clear audit trail from informal theory statements to formal proof
artifacts.

\subsection{Full Correspondence Tables}

Detailed chapter-wise mappings are provided in Appendix~F. A compact summary mapping is included in Appendix~A for audit continuity.

\subsection{Coverage Summary}

The correspondence tables are designed so that each named theory claim has at
least one Lean anchor theorem, and each novel theorem carries an explicit
annotation indicating that it was not a verbatim carry-over from the original
theory manuscript.
